\section{Discussion}
We interpret the results contribution by contribution: the design choices and scope of the C1 dataset, the classifier-design rationale and cross-distribution behavior of C2, and the architectural rationale, abstraction-level findings, and workflow-fit signals of C3.


\subsection*{C1: Dataset Quality and Scope}\label{sec:discussion-c1}

Wu et al.~\cite{wu2022data-quality} document substantial label noise in community-assigned security fields across the widely used Jira-based benchmarks, and subsequent classifier work~\cite{alqahtani2024fasttext, franca2025gpts, soltaniani2026evaluating-llms} inherits that noise even after Wu's cleaning pass.
C1 trades scale for annotation rigor: $288$ items fully multi-rated by $9$ security experts produce $972$ independent assessments at Fleiss' $\kappa = 0.787$.
The design target is to complement Wu et al.'s dataset with a high-precision training signal whose labels were assigned by practitioners whose day-to-day role is to decide exactly the kind of question the annotation asks.
Wu et al.'s datasets remain the larger-scale evaluation. C1 is the clean training supervision.

The $13.5\%$ of items with at least one dissenting assessor, despite substantial inter-rater agreement overall, confirm that security relevance is not a sharp classification boundary.
A natural generalizability question for any expert-annotated resource is whether a different expert pool would converge on similar labels.
The annotators share a common organizational context, but that context is broad due to the size of the industrial partner, which operates across multiple highly regulated domains and is therefore subject to a wide cross-section of security standards and regulations.
This implies that the practitioners' day-to-day decisions span much of the compliance landscape that other security-regulated organizations also face, making their experience representative of many industries rather than narrowly organization-specific.
The working definition of \texttt{security-relevant} that they apply is in turn shaped by the same public standards (e.g., IEC 62443, CIS Benchmarks) that those organizations follow, so it is not idiosyncratic to a single industry.

The source corpus is intentionally narrow: two Atlassian enterprise projects whose item language resembles the industrial backlogs targeted by C3.
Datasets whose item descriptions are tersely phrased or centered on consumer-facing bug reports sit at a different item granularity and are better served by the Wu family.
The $80$ security-related items also bound what future work can extract from C1 alone: frozen-encoder classifiers trained as in C2 reach the ceiling of this signal comfortably, but multi-class Firesmith-category models~\cite{firesmith2004specifying} would require augmentation or a second annotation round.

\paragraph{Using C1 in practice.}
Beyond its role as training supervision for C2, we see three concrete uses of the released dataset for industry practitioners.
First, as expert-labeled seed data to bootstrap an organization-internal security-relevance classifier before internal labels exist, following the C2 procedure of a frozen encoder with a lightweight head.
Second, as an audit set. Organizations already operating a triage classifier or LLM-based filter can evaluate it against expert consensus labels rather than community-assigned ones.
Third, as realistic material for security training. The items are industry-style backlog texts with expert consensus security-relevance labels, which teams can use to align their own working definition of security relevance.

\subsection*{C2: Classifier Design Decisions}\label{sec:discussion-c2}

With $288$ annotated items, full fine-tuning of a contemporary transformer risks overfitting to the two source projects rather than learning a transferable notion of security-relevance.
Linear-head coefficients are also inspectable: security-filtering decisions are reviewed internally, and an auditable decision path matters for trust and for error analysis.
F2 as the selection criterion encodes the corresponding cost asymmetry fairly, since a missed security requirement is more expensive than a false alarm that C3 subsequently filters through its semantic stage.


\subsubsection*{Generalizability evaluation}\label{sec:discussion-c2-generalizability}

The generalizability results should be interpreted as evidence of transferability.
The deployed C2 model is evaluated fully zero-shot across domains: trained only on C1 and then applied unchanged to the Wu et al.\ benchmarks, whereas most published baselines are evaluated in-distribution, with training and test data drawn from the same project.
Matching or approaching those baselines is therefore a comparatively demanding target rather than a neutral head-to-head comparison.

The comparison against prompted LLMs should also be read cautiously. Although they are evaluated in a nominally zero-shot setting, their pretraining data may plausibly already contain parts of these widely used benchmark projects, potentially resulting in overly optimistic results~\cite{deng2024llm-contamination, baltes2026guidelinesempiricalstudiessoftware}. We therefore treat competitive performance on Wu et al.\ as a positive result showing that C2 has learned a transferable notion of \texttt{security-relevant}.

The reported superior performance of Gemini by Soltaniani et al.~\cite{soltaniani2026evaluating-llms}, while impressive, cannot be perceived as a stable baseline given recent insights on the non-reproducibility of research performed with commercial LLMs~\cite{angermeirReproducibility}.

The tight clustering of in-distribution F2 across eight sentence-encoder-based combinations (F2 $\in [0.625, 0.824]$) suggests that the task is well-posed and recoverable by any competitive text representation. The bottleneck here is label quality.
DistilRoBERTa + LogReg (L2) has the highest held-out F2 overall, but was not treated as a deployment candidate because on a two-project training resource its $\ell_2$ head is more exposed to project-specific overfitting than the sparser elastic-net alternatives.
Among the transfer-oriented deployment candidates, MPNet + LogReg (Elastic) has the highest held-out F2 and was therefore selected for integration into the C3 retrieval system.
The cross-distribution comparison uses mean G-measure as its primary transfer metric and across-project standard deviation $\sigma$ as a secondary stability metric.
Under that criterion, MPNet + elastic-net logistic regression remains a strong deployment choice, but the stability comparison also surfaces a real trade-off: MiniLM-L12 + elastic-net logistic regression achieves essentially identical mean G-measure ($0.650$ vs.\ $0.649$) with half the cross-project standard deviation ($\sigma = 0.046$ vs.\ $0.102$).
The fact that both variants are competitive with the strongest within-project baselines on stability, and that fine-tuned small LLMs are considerably less consistent despite per-project training supervision, supports the interpretation that the classifier's stability derives from the transferability of the learned representation.
In deployments where cross-domain consistency is the primary concern, MiniLM-L12 + LogReg (Elastic) would therefore be the preferred variant.

Probability calibration is deliberately deferred: C2 is used downstream as a hard filter into C3, where calibration does not affect retrieval-pipeline behavior.
This would matter only if C2 output fed a downstream ranking, which is not the case in the current deployment.

\subsection*{C3: Security Requirement Enrichment}\label{sec:discussion-c3}

We discuss C3 along its design rationale, the practical value to developer workflows observed in the preliminary evaluation, and four findings that emerged in individual walkthroughs.

\subsubsection*{Design rationale}

\emph{Two-stage C2--C3 system.}
Separating the high-recall relevance filter from the precision-oriented retrieval stage matches the asymmetric cost of each stage and lets each subsystem be tuned independently: C2 maximizes recall of security-relevant items, C3 maximizes precision of clause grounding.
The LLM workload is therefore proportional to the number of security-relevant items, not to the full backlog size.

\emph{RAG over pure LLM IR.}
The bi-encoder shortlist imposes a document-grounded candidate set that the LLM filter cannot hallucinate out of, while the LLM filter rejects topically adjacent but non-applicable candidates that dense similarity alone cannot distinguish.

\emph{Requirement extraction preprocessing.}
Introduced after we observed a systematic failure mode where multi-concern backlog items diluted the retrieval signal of any single concern.
Decomposing into atomic requirements before retrieval recovers concerns that the undecomposed retrieval pipeline misses.

\emph{Open-weights LLMs.}
Using on-premises open-weights models for both extraction and filtering satisfies two deployment constraints simultaneously: backlog contents cannot be sent to third-party inference services, and pinning a specific model version prevents silent behavior drift as commercial models are updated. More of the benefits of open LLMs can be found in the LLM guidelines by Baltes et al.~\cite{baltes2026guidelinesempiricalstudiessoftware}.

\subsubsection*{Practical value}

The tool surfaces implicit security requirements that would otherwise require manual expert review, delivering security requirements at the point of planning rather than as an after-the-fact compliance check, potentially addressing the lead-time burden identified by Moyón et al.~\cite{moyon2024challenges}.
In the preliminary evaluation, both practitioners responded positively to this framing. The tool was perceived as useful at the item level, provided it integrates into existing workflows rather than introducing a parallel interface.

\subsubsection*{Abstraction-level mismatch}\label{sec:discussion-c3-abstraction}

Across both interview sessions, the granularity gap between security requirement documents and backlog items appeared to be the strongest predictor of retrieval usefulness.
Rule-level documents, such as organizational security policies and implementation-level benchmarks such as CIS Benchmarks, operate below the capability abstraction of standards such as IEC 62443-3-3~\cite{iec2013iec62443-3-3}, and yielded actionable matches in both cases.
However, the within-shortlist variance of CIS Benchmark retrievals was higher, consistent with their narrower implementation scope making them precise for specific infrastructure items but less broadly applicable across item types.
Higher-abstraction capability standards, such as IEC 62443-3-3, produced retrievals that were topically correct but rarely directly actionable per item. The semantic gap between such control and a backlog item describing a specific feature is too large for retrieval to bridge reliably.
If confirmed by the upcoming complete evaluation, this finding has a direct implication for deployment.
The retrieval pipeline should be paired with documents whose security requirements granularity matches the typical specificity of the target backlog.

A secondary observation is that retrieval shortlist quality appears to track backlog-item granularity itself. Coarse epics retrieve many weakly relevant security requirements, while atomic tasks retrieve fewer but more strongly relevant ones.
This suggests a potential secondary use of C2\,+\,C3 output as a proxy signal for backlog-structuring quality, though we do not validate this claim in this paper.

\subsubsection*{Workflow fit}

Both practitioners independently identified in-backlog-system integration as the primary precondition for practical adoption.
Standalone interfaces were perceived to add tool-fatigue. Retrieval output surfaced as issue comments or side-panel content can be consumed at the point where planning actually happens.
This is considered the most consistent qualitative finding of the interview study.

Security requirement verbosity was the dominant usability concern. Retrieved clauses from both internal policies and public benchmarks are often too long for inline reading.
A natural next engineering step is to use the same LLM stage to generate a short, concrete action-item summary per retained security requirement, with a link back to the authoritative text for verification, an addition that has been implemented but must be evaluated.
Perceived usefulness also appeared to correlate with source-item quality. Items with structured descriptions produced higher-rated retrievals than items with terse, unstructured titles.


\subsubsection*{Missing items and concerns}

The interviews produced few systematic gap reports.
The one recurring finding was that the source corpus occasionally omits cross-cutting requirements. Backwards-compatibility obligations flagged in one session are not expressed as security clauses in any of the evaluated standards, and the tool cannot surface what is not in the corpus.
The extractor continues to have difficulties with multi-concern items, which is why further multi-concern decomposition is an open direction.
This limitation may also be mitigated by establishing guidelines for backlog item descriptions, since more detailed items may provide better source material for retrieving relevant security requirements.

\subsubsection*{Generalizability}
The C2 classifier is trained entirely on open-source TAWOS data and carries no organization-specific knowledge, so adopters applying C3 to their own internal security policies need no re-annotation or retraining.
The strongest generalizability claim for C3, that the retrieval pipeline performs consistently across standards and organizations, is part of future work.
The CIS Benchmark evaluation partially addresses this by demonstrating consistent pipeline behavior on a public security requirements document against the same items used for the internal-policy evaluation, but a second independent industrial corpus would be needed to convert the architectural document-agnosticism claim into a measured result.
